\vspace*{2cm}
\begin{center}
\Large\textbf{APPENDIX}
\end{center}
\vspace{1.5cm}

\section*{A \quad PhoBERT Sentiment Model Details}

Our sentiment classification pipeline leverages PhoBERT, a Vietnamese BERT variant pretrained on Vietnamese Wikipedia and news corpora. We fine-tune a pretrained PhoBERT-base model (12 layers, 768 hidden dimensions) on a CafeF article subsample labeled as positive (positive returns/optimistic tone), negative (negative returns/pessimistic tone), and neutral using the transformers library. The fine-tuning uses a standard cross-entropy loss, AdamW optimizer (learning rate $1 \times 10^{-5}$), and 3 epochs over a labeled subset of 2,000 articles. We extract sentiments at the article level by tokenizing the title and lead paragraph (first 512 tokens), applying the fine-tuned model, and assigning a sentiment class (positive, negative, neutral) based on the highest logit. The sentiment classes are then mapped to numerical scores: positive = \(+1\), neutral = \(0\), negative = \(-1\). Daily sentiment aggregates compute the mean score across all articles published that day, producing a daily sentiment index ranging from $-1$ to $+1$ with high positive skew (84\% positive articles in 2015--2024 CafeF corpus).

\section*{B \quad Data Preprocessing and Trading-Day Alignment Rules}

Raw CafeF articles and OHLC prices were aligned to a common daily frequency using the following protocol: (1) scrape article timestamps (publication date and time) from CafeF; (2) categorize articles by publication hour—articles published before 14:45 (market close) are attributed to day $t$'s sentiment, articles after 14:45 are attributed to day $t+1$; (3) exclude weekends and Vietnamese holidays (New Year, Lunar New Year, National Day, etc.) from the trading calendar; (4) compute log-returns as $r_t = \log(C_t / C_{t-1})$, where $C_t$ is the closing price on trading day $t$; (5) winsorize returns at 2.5\% tails to remove outliers; (6) align each trading day to its corresponding sentiment score, lagged LSTM features, and GARCH residuals. The resulting dataset spans 2,328 trading days (2015--2024) with complete price, sentiment, and feature vectors for all training and test periods.

\section*{C \quad Extended Robustness Results Across Alternative Specifications}

To validate the stability of our findings beyond the baseline configuration, we evaluate the hybrid model under five alternative specifications: (1) \textbf{Alternative Sentiment Thresholds:} Instead of mapping sentiment into $\{-1, 0, +1\}$, we test $\{-1, -0.5, 0, +0.5, +1\}$ (finer granularity) and $\{-1, +1\}$ (binary). Results remain robust with RMSE improvements of 14.2\%--14.6\% across thresholds, confirming that the hybrid architecture benefits are not driven by a specific discretization choice. (2) \textbf{Garman-Klass Volatility Target:} Replacing the realized volatility target (based on log high/low) with the Garman-Klass estimator yields a 16.05\% RMSE improvement, indicating the model generalizes to alternative volatility metrics. (3) \textbf{Feature Ablation:} Removing sentiment features from the LSTM input reduces RMSE improvement to 14.26\%, showing that the nonlinear architecture (not sentiment alone) drives 99\% of gains. (4) \textbf{Linear GARCH-X Baseline:} Direct embedding of sentiment into the variance equation via GARCH(1,1)-X fails spectacularly ($t_{DM} = -6.339$, RMSE degradation of 5.8\%), confirming that nonlinearity is essential. (5) \textbf{Expanding-Window Estimation:} Using rolling 500-day windows with expanding start dates yields 13.9\% RMSE improvement, confirming robustness under non-stationary retraining protocols.

\section*{D \quad CafeF Corpus Statistics and Data Quality Metrics}

The CafeF news corpus spans 2,328 trading days (January 2015--December 2024) and comprises 86,179 unique articles. Temporal coverage is uneven: 2015--2019 average 10.2 articles per trading day; 2020--2024 average 16.8 articles per day, reflecting increased financial news production in recent years. Sentiment distribution exhibits a strong positive skew: 84.1\% positive articles, 13.2\% neutral, 2.7\% negative. This imbalance persists across subperiods and mirrors retail-investor optimism bias in Vietnam's emerging market. Article length (title + lead) averages 78 tokens, with 95th percentile at 256 tokens. Data quality checks reveal 0.3\% duplicate articles (removed), 1.2\% articles with malformed timestamps (imputed to nearest trading day), and 0.8\% articles with blank sentiment text (assigned neutral class). Missing trading days (holidays) are handled by forward-filling sentiment from the prior trading day; empirically, this affects 2.1\% of observations and does not materially affect results (sensitivity analysis on file upon request).

\section*{E \quad LSTM Architecture and Hyperparameter Selection}

The LSTM residual-correction model processes 15-day lagged sequences of GARCH residuals and sentiment features into 1-step-ahead volatility forecasts. The architecture consists of an LSTM layer (128 hidden units, dropout = 0.2) followed by two fully connected layers (64 units with ReLU activation, dropout = 0.2; 1 output unit with linear activation). Input features include: (i) lagged GARCH residuals ($r_{t-1}, \ldots, r_{t-15}$), (ii) lagged sentiment scores ($s_{t-1}, \ldots, s_{t-15}$), (iii) news intensity (daily article count, clipped at 99th percentile), and (iv) volatility regime indicator (realized volatility binned into quintiles). Hyperparameters were selected via grid search over the 2015--2021 training window: LSTM hidden units $\in \{64, 128, 256\}$, dropout $\in \{0.1, 0.2, 0.3\}$, fully-connected layer width $\in \{32, 64, 128\}$, learning rate $\in \{1 \times 10^{-4}, 5 \times 10^{-4}, 1 \times 10^{-3}\}$. Optimal selection (128 LSTM units, 0.2 dropout, 64 dense units, $5 \times 10^{-4}$ learning rate) was determined by validation RMSE on a held-out 2021 window. Early stopping monitored validation loss with patience = 10 epochs. Training used Adam optimizer and batch size = 32 over 100 epochs, achieving convergence by epoch 67 on average.

